% Motivationsschreiben DI Alexander Sommer - Bewerbung FMA Österreichische Finanzmarktaufsicht
% Position: Teamleitung Strategische Steuerung und Innovation (Bewerbungsfrist 20.11.2026)
% Compiler: Tectonic / pdfLaTeX
\documentclass[10.5pt]{article}
\usepackage[a4paper,top=16mm,bottom=15mm,left=21mm,right=21mm]{geometry}
\usepackage[T1]{fontenc}
\usepackage[utf8]{inputenc}
\usepackage[ngerman]{babel}
\IfFileExists{sourcesanspro.sty}%
  {\usepackage[default,semibold]{sourcesanspro}}%
  {\usepackage[scaled=0.92]{helvet}\renewcommand{\familydefault}{\sfdefault}}
\usepackage{microtype}
\usepackage{xcolor}
\usepackage[hidelinks]{hyperref}

\definecolor{accent}{HTML}{1F3A5F}
\definecolor{brandblue}{HTML}{1E3A8A}
\definecolor{muted}{HTML}{5A6472}

\hypersetup{
  pdftitle={Motivationsschreiben DI Alexander Sommer - FMA},
  pdfauthor={DI Alexander Sommer}
}

\setlength{\parindent}{0pt}
\setlength{\parskip}{4.5pt}
\linespread{1.04}
\pagestyle{empty}

\begin{document}

% ---------------------------------------------------------------- Absender / Empfänger
\begin{minipage}[t]{0.56\textwidth}
  \textbf{\color{accent}\large DI Alexander Sommer}\\[1pt]
  Gregorygasse 12/6/2 · 1230 Wien\\[1pt]
  Tel.: +43 681 102 104 17 · \href{mailto:alexander.sommer@gmail.com}{alexander.sommer@gmail.com}\\[1pt]
  \href{https://alexander-sommer.vercel.app}{alexander-sommer.vercel.app}
\end{minipage}\hfill
\begin{minipage}[t]{0.42\textwidth}\raggedleft
  \textbf{Österreichische Finanzmarktaufsicht (FMA)}\\
  Personalmanagement\\
  Otto-Wagner-Platz 5 · 1090 Wien
\end{minipage}

\vspace{8pt}
\hfill Wien, 8. November 2026

\vspace{2pt}
{\large\bfseries\color{brandblue}Bewerbung für die Teamleitung Strategische Steuerung und Innovation}

\vspace{4pt}
Sehr geehrte Damen und Herren,

„Kompetenz schafft Vertrauen“ -- dieser Leitsatz der FMA beschreibt für mich sehr treffend, worauf es in der Strategie- und Veränderungsarbeit ankommt: Orientierung entsteht dort, wo Ziele nachvollziehbar sind, Vorhaben konsequent umgesetzt werden und Führungskräfte sich auf eine verlässliche Steuerung stützen können. An genau dieser Schnittstelle zwischen Strategie, Organisation und Umsetzung arbeite ich seit mehr als 15 Jahren -- bei der OMV, bei Wiener Wohnen und derzeit als Teamleiter Digital Workplace Solutions bei der Austrian Power Grid AG (APG). Nach vielen Jahren in der kritischen Infrastruktur möchte ich diese Erfahrung nun gerne in einer Organisation einbringen, deren öffentlicher Auftrag -- die Stabilität des Finanzmarkts und das Vertrauen in ihn -- unmittelbare gesellschaftliche Relevanz hat.

Was ich für die ausgeschriebene Aufgabe mitbringe, lässt sich entlang Ihrer Schwerpunkte zusammenfassen:

\begin{itemize}\setlength{\itemsep}{1.5pt}\setlength{\parskip}{0pt}
  \item \textbf{Strategie- und Zielentwicklung:} Als Leiter Digitalisierung \& Innovation bei Wiener Wohnen Kundenservice verantwortete ich die Entwicklung und Umsetzung der Digitalisierungs- und Innovationsstrategie eines kommunalen Unternehmens und beriet die Geschäftsführung in strategischen Fragen. Dort habe ich gelernt, dass eine Strategie im öffentlichen Auftrag nur dann trägt, wenn sie in überprüfbare Ziele übersetzt und von den Führungskräften aktiv mitgetragen wird.

  \item \textbf{Projektportfolio \& Umsetzungsstärke:} Bei der OMV leitete ich mehr als 40 internationale IT-Projekte und habe konzernweite Einsparpotenziale von bis zu 4\,Mio.\,EUR pro Jahr strukturiert identifiziert. Bei der APG verantwortete ich die Steuerung und Harmonisierung eines Portfolios von mehr als 340 Applikationen. Priorisieren, Abhängigkeiten sichtbar machen und Fortschritt ehrlich berichten -- das ist für mich der Kern eines wirksamen Portfoliomanagements.

  \item \textbf{Organisationsentwicklung \& Change:} Als Kernmitglied des konzernweiten Veränderungsprogramms \emph{Change4Growth} habe ich mein damaliges Team bewusst neu aufgestellt und näher an die strategische Entscheidungsebene gerückt; die daraus entstandene Schwesterstruktur steuere ich bis heute in einer Matrixfunktion mit. Veränderung gelingt aus meiner Erfahrung dann, wenn Betroffene früh beteiligt werden und den Sinn hinter einer Maßnahme verstehen.

  \item \textbf{Innovation mit Augenmaß:} Von der Innovationsjury und dem Start-up-Ökosystem der OMV über meine Rolle als \emph{Innovation Peer} bei der APG bis zum Programm \emph{Applied Agentic AI} (MIT, 2025) und sieben Jahren als Gartner Seat Holder: Neue Technologien beurteile ich lieber selbst, als mich auf Trends zu verlassen. Als Miterfinder eines vom Österreichischen Patentamt erteilten, FFG-geförderten Patents (AT\,519762) kenne ich zudem den Weg von der ersten Idee bis zur geschützten, umsetzbaren Innovation aus eigener Erfahrung. Innovation verstehe ich dabei als kontinuierliche Verbesserung, die sich an ihrem Nutzen für Organisation und Menschen messen lassen muss.
\end{itemize}

Führung bedeutet für mich, Orientierung zu geben, Verantwortung zu übertragen und eine offene, leistungsorientierte Lern- und Feedbackkultur vorzuleben. Diese Haltung wurde 2023 auch extern bestätigt: In einer Management-Potenzialeinschätzung (PEP, IMDE) lag mein Führungsprofil im oberen Dezil vergleichbarer Führungskräfte. Besonders freut mich, dass Mitarbeitende, die unter meiner Führung begonnen haben, heute selbst Verantwortung tragen -- etwa Bernhard Mittermann, inzwischen CIO bei Rexel. Gerne nenne ich Ihnen auf Wunsch Referenzen.

Die Welt der Finanzmärkte ist mir dabei nicht fremd: Als Mitgründer von 42Investors habe ich regelbasierte Handelssysteme für NYSE und NASDAQ mitentwickelt, mit Schwerpunkt auf Risikoparametern und Systemstabilität im Live-Betrieb (\emph{Wall Street Journal}, 10.~August 2015). Die aufsichtsrechtliche Fachtiefe Ihrer Expertinnen und Experten ersetzt das selbstverständlich nicht -- mein Beitrag läge vielmehr darin, diese Expertise durch klare Ziele, eine verlässliche Steuerung und gute Zusammenarbeit über Bereichsgrenzen hinweg bestmöglich wirksam werden zu lassen.

Über die Gelegenheit, Sie in einem persönlichen Gespräch kennenzulernen, würde ich mich sehr freuen.

\vspace{6pt}
Mit freundlichen Grüßen\\[16pt]
\textbf{DI Alexander Sommer}

\end{document}
